\documentclass[12pt,a4paper]{article}

% -------------------- Packages --------------------
\usepackage[utf8]{inputenc}
\usepackage[T1]{fontenc}
\usepackage{mathpazo} 
\usepackage[margin=1in]{geometry}
\usepackage{xcolor}
\usepackage{fancyhdr}
\usepackage{titlesec}
\usepackage{listings}
\usepackage[hidelinks]{hyperref}

% -------------------- Styling --------------------
\definecolor{LibraryBlue}{RGB}{0, 51, 102}
\definecolor{CodeGray}{RGB}{240, 240, 240}

% Headings
\titleformat{\section}{\color{LibraryBlue}\Large\bfseries}{\thesection}{1em}{}[\titlerule]
\titleformat{\subsection}{\color{LibraryBlue}\large\bfseries}{\thesubsection}{1em}{}

% Code Blocks
\lstset{
    backgroundcolor=\color{CodeGray},
    basicstyle=\ttfamily\small,
    frame=single,
    breaklines=true,
    columns=fullflexible
}

% Header/Footer
\pagestyle{fancy}
\fancyhf{}
\fancyhead[L]{\textit{Deployment Documentation}}
\fancyhead[R]{\thepage}
\renewcommand{\headrulewidth}{0.4pt}

\begin{document}

\section{Deployment \& DevOps}

This chapter details the technical infrastructure, build processes, and continuous integration pipelines established for the Library Mobile Application.

\subsection{Development Environment Configuration}
The application is built using the **Expo** framework on top of **React Native**. Setting up a consistent development environment is critical for reproducibility.

\subsubsection{Prerequisites}
\begin{itemize}
    \item \textbf{Node.js}: Runtime environment (v20+ recommended).
    \item \textbf{Expo CLI}: Command-line interface for running and managing Expo projects.
    \item \textbf{EAS CLI}: Expo Application Services tool for cloud builds.
    \item \textbf{Firebase}: Helper tools for backend management.
\end{itemize}

\subsection{Build Architecture (EAS)}
We utilize **Expo Application Services (EAS)** for generating production-ready binaries. The configuration is defined in \texttt{eas.json}, specifying three distinct environments:

\begin{enumerate}
    \item \textbf{Development}: Generates a development client for rapid testing on physical devices.
    \item \textbf{Preview}: Produces a standalone APK for internal distribution and QA testing.
    \item \textbf{Production}: optimized builds signing for Google Play Store submission.
\end{enumerate}

\subsubsection{Manual Build Command}
To trigger a manual build for Android, the following command is executed:
\begin{lstlisting}
eas build --platform android --profile preview
\end{lstlisting}

\subsection{Continuous Integration (CI/CD)}
To ensure code quality and automate delivery, we have implemented a CI/CD pipeline using **GitHub Actions**.

\subsubsection{Workflow Definition}
The workflow is defined in \texttt{.github/workflows/build.yml} and is triggered on:
\begin{itemize}
    \item Pushes to the \texttt{main} branch.
    \item Pull Requests targeting the \texttt{main} branch.
\end{itemize}

\subsubsection{Pipeline Steps}
The automated pipeline executes the following sequence:
\begin{enumerate}
    \item \textbf{Checkout}: Retrieves the latest code from the repository.
    \item \textbf{Setup}: Initializes the Node.js environment.
    \item \textbf{Dependency Installation}: Runs \texttt{npm ci} for a clean install.
    \item \textbf{Authentication}: Authenticates with Expo using the \texttt{EXPO\_TOKEN} secret.
    \item \textbf{Build}: Executes the non-interactive EAS build process.
    \item \textbf{Artifact Upload}: Stores the resulting APK as a downloadable artifact for review.
\end{enumerate}

\subsection{Backend Integration}
The application relies on **Firebase** for backend services.
\begin{itemize}
    \item \textbf{Authentication}: Manages user sessions and academic identities.
    \item \textbf{Firestore}: NoSQL database for storing user profiles, book catalogs, and borrowing history.
\end{itemize}
API keys are securely managed via environment configurations files (`firebaseConfig.js` and `gemini.js`).

\end{document}
